\subsection*{MT-S504. The exact relative central condition in the $S_5$ transfer}

Retain the $S_5$ extension and local hypotheses of MT-S501:
complex conjugation and the unramified Frobenius at $5$ both
have cycle type $(12)(34)$. Write $E=L^{A_5}$, let $F$
be a degree-six subfield, and let $H=EF$. Let $\eta$ be the
quadratic character of $E/\mathbb Q$. Use Calegari's labeling
of the five-dimensional character $\rho_5$, whose values on
the classes $1,2A,2B,3A,4A,5A,6A$ are
$(5,1,1,-1,-1,0,1)$. Then the all-prime Artin factorization is
\[
       \zeta_H(s)=\zeta_F(s)L(s,\eta)L(s,\rho_5).
 \tag{S504.1}
\]
Under the conditional input $M_H$, the strictly more focused
sufficient condition
\[
 \operatorname{ord}_{1/2}\zeta_H
       \le\operatorname{ord}_{1/2}\zeta_F
                     +\operatorname{ord}_{1/2}L(s,\eta)
 \tag{S504.2}
\]
implies the same conclusion as MT-S502: every nontrivial
irreducible Artin factor is entire and satisfies GRH, and
the factors are automorphic by Calegari's argument.
In fact equality then holds in (S504.2). The hypothesis
allows zeros in the two other factors; it does not require
$\zeta_H(1/2)\ne0$.

\begin{proof}
Put $m_5=\operatorname{ord}_{1/2}L(s,\rho_5)$, interpreted
as the meromorphic order, so negative values are poles.
Equation (S504.1) shows that (S504.2) is exactly $m_5\le0$.
The factorization and the following precise classical inputs
are Calegari, \emph{The Artin conjecture for some
$S_5$-extensions}, \texttt{arXiv:1112.1152v1}: the comment
after Theorem~1.2, p.~2; Lemmas~4.8--4.9, pp.~10--11;
and the final proof of Theorem~1.2, p.~11. Their relevant
proofs have been read in full. The automorphic construction
in preceding sections remains a cited classical input.

That construction supplies a cuspidal automorphic
$\varpi$ on $\mathrm{GL}_5/\mathbb Q$, an exceptional set
$\Omega$ of primes, and a meromorphic comparison
\[
 \Lambda(s,\varpi)=\Lambda(s,\rho_5)\mu(s),\qquad
 \mu(s)=\mu_B(s)
       \prod_{p\in\Omega}
             \left(\frac{1+p^{-s}}{1-p^{-s}}\right)^2.
 \tag{S504.3}
\]
Here $\mu_B$ is holomorphic and nonzero on $(0,1)$.
Lemma~4.8 says that if $\Omega$ is infinite, $\mu$ has
a pole at a real point in $(0,1)$. Lemma~4.9 says that
finite $\Omega$ forces exact local equality everywhere,
including ramified and archimedean factors. The completed
automorphic numerator in (S504.3) is holomorphic on this
interval.

Brauer induction and $M_H$ imply that $L(s,\rho_5)$ is
holomorphic and nonzero at every real point of $(0,1)$
other than $1/2$: every Hecke factor in its signed
factorization has that property. Its archimedean factors
are finite and nonzero on the interval. Thus a real pole
of $\mu$ could only occur at $1/2$ and would require
$m_5>0$. Our hypothesis $m_5\le0$ rules that out;
a pole of the denominator would instead give a zero
of the comparison. Lemma~4.8 makes $\Omega$ finite,
and Lemma~4.9 proves $\Lambda(s,\varpi)=\Lambda(s,\rho_5)$.
The source's final argument also gives the exact match
for the quadratic twist $\rho_5\otimes\eta$.

The other irreducible dimensions $4$ and $6$ are already
automorphic by Theorem~1.1, as recorded in MT-S501;
the linear factors are classical Dirichlet functions.
This proves all the automorphy and holomorphy assertions.
Now holomorphy makes $m_5\ge0$, so $m_5=0$ and equality
holds in (S504.2). Finally signed Brauer factorization
and $M_H$ locate the nontrivial zeros of every entire
factor on the critical line, exactly as in MT-ART07.
\end{proof}

All $S_5$ irreducibles are realizable over $\mathbb R$;
Fr\"ohlich--Queyrut's real-representation theorem gives
root number $+1$. The meromorphic central-order parity
calculation of MT-ART11 therefore makes $m_5$ even.
Consequently failure of the sufficient condition leaves
only the gate
\[
 \operatorname{ord}_{1/2}\zeta_H
       -\operatorname{ord}_{1/2}\zeta_F
       -\operatorname{ord}_{1/2}L(s,\eta)
                       =m_5\ge2.
 \tag{S504.4}
\]
This is a necessary condition for this route to remain
unsettled, not a claim that a positive central order
prevents automorphy. It removes base-factor central zeros
from the test exactly; no independence or noncoincidence
of zeros is assumed.

Geometrically the virtual permutation identity in
(S504.1) isolates a single rank-five Artin motive from
the arithmetic finite covers. It sharpens the
MT-S502 one-value gate in the cote~1/SGA5 trace-divisor
programme, and the $S_5$ portion of Q-ARTIN-POS.
It leaves the actual relative central order as an
arithmetic input; neither designated GRH theorem
evaluates it.
